\chapter{Conception et architecture}
\label{ch:archi}

\section{Architecture globale}
POSTIE suit une architecture trois-tiers strictement découplée
(\cref{fig:archi})~: un frontend Next.js servi par Vercel, un backend
FastAPI conteneurisé sur Render, et un vector store ChromaDB
embarqué dans le conteneur backend. Les modèles de langage GPT-4o
et Mistral Large sont consommés via leurs API HTTPS respectives.

\begin{figure}[H]\centering
\begin{tikzpicture}[
  node distance=1.2cm and 1.4cm,
  block/.style={rectangle,draw=postieblue,thick,rounded corners=4pt,
                fill=lightgray,minimum width=3.4cm,minimum height=1.1cm,
                align=center,font=\small},
  ext/.style  ={rectangle,draw=accentviolet,thick,rounded corners=4pt,
                fill=accentviolet!10,minimum width=3.4cm,minimum height=1.1cm,
                align=center,font=\small},
  arr/.style  ={-{Stealth[length=3mm]},thick,postieblue}
]

\node[block] (front) {\textbf{Frontend}\\Next.js 15 + React 19\\Tailwind};
\node[block,right=of front] (back) {\textbf{Backend API}\\FastAPI + LangChain\\Python 3.11};
\node[block,right=of back]  (vec)  {\textbf{Vector Store}\\ChromaDB};

\node[ext,below=of front]   (vercel){Vercel (CDN)};
\node[ext,below=of back]    (render){Render (Docker)};
\node[ext,below=of vec]     (llm)   {OpenAI GPT-4o\\Mistral Large};

\draw[arr] (front)  -- node[above,font=\scriptsize]{HTTPS / JSON} (back);
\draw[arr] (back)   -- node[above,font=\scriptsize]{MMR k=6}      (vec);
\draw[arr] (back)   -- node[right,font=\scriptsize]{LLM API}      (llm);
\draw[arr] (front)  -- (vercel);
\draw[arr] (back)   -- (render);

\begin{scope}[on background layer]
  \node[draw=postieblue!30,dashed,rounded corners,inner sep=8pt,
        fit=(front)(back)(vec),
        label=above:{\small\bfseries Application POSTIE}] {};
\end{scope}

\end{tikzpicture}

\caption{Architecture globale de POSTIE.}
\label{fig:archi}
\end{figure}

\section{Stack technique et justification}
\Cref{tab:stack} récapitule les choix technologiques majeurs et leur
justification.

\begin{table}[H]\centering\small
\caption{Stack technique et justification des choix.}
\label{tab:stack}
\begin{tabularx}{\linewidth}{llX}
\toprule
\textbf{Couche}      & \textbf{Technologie}     & \textbf{Justification}\\
\midrule
Frontend             & Next.js 15, React 19, TS & SSR, écosystème mature, typage strict\\
Backend              & FastAPI 0.115            & Async natif, typage Pydantic, OpenAPI auto\\
Orchestration LLM    & LangChain 0.3            & Chaînes RAG prêtes, mémoire conversationnelle\\
Vector store         & ChromaDB 0.5             & Embarqué, persistant, simple à déployer\\
Embeddings           & text-embedding-3-small   & Rapport coût/qualité optimal\\
LLM principal        & GPT-4o                   & Raisonnement, génération longue\\
LLM secondaire       & Mistral Large            & Scoring structuré, sortie JSON\\
Déploiement front    & Vercel                   & Auto-deploy GitHub, CDN, gratuit\\
Déploiement back     & Render                   & Docker, free tier, disque persistant\\
Contrôle de version  & Git + GitHub             & Workflow trunk-based, CI/CD intégré\\
\bottomrule
\end{tabularx}
\end{table}

\section{Schéma de données}
Les features du backlog sont représentées par le modèle Pydantic
\texttt{Feature} (\cref{lst:feature}) qui agrège l'ensemble des
critères nécessaires aux sept algorithmes de priorisation.

\begin{lstlisting}[caption={Modèle Pydantic Feature (extrait).},label={lst:feature}]
class Feature(BaseModel):
    id: str
    name: str
    description: str
    epic: str
    tags: list[str]
    # Critères RICE
    reach: float
    impact: float
    confidence: float
    effort: float
    # Critères WSJF
    business_value: float
    time_criticality: float
    risk_reduction: float
    job_size: float
    # Critères ICE / Value-Effort
    ease: float
    # Classifications
    kano_category: Literal["basic", "performance", "delighter"]
    moscow: Literal["must", "should", "could", "wont"]
\end{lstlisting}

\section{Sécurité et déploiement}
\begin{itemize}
  \item Clés API injectées uniquement par variables d'environnement
        (jamais commitées dans le dépôt).
  \item \gls{cors} restreint à une whitelist explicite plus une regex
        \texttt{*.vercel.app} pour les previews et
        \texttt{localhost:*} pour le développement local.
  \item Audit \texttt{npm audit} et \texttt{pip-audit} avant chaque
        release.
  \item Sanitization de l'historique Git après détection d'un
        \emph{leak} accidentel d'une clé OpenAI (sprint 6).
  \item Mise à jour Next.js 15.0.3~$\rightarrow$~15.5.18 pour
        corriger la CVE-2025-66478 bloquant le déploiement Vercel.
\end{itemize}
